\subsection{13. FDTD: Solving 1+1D Delay PDE in Parallel}
\textbf{OSTI:} \texttt{1475143}\quad\textbf{Authors:} Zheng, Goldschmidt, Fan (2018)\quad\textbf{Domain:} Quantum Optics / Numerical PDE\\
\scorebadge{Coverage}{7}\quad\scorebadge{Agreement}{7}\quad\textbf{Total:} 14/20\par\smallskip
\noindent\textbf{Coverage rationale.} Core FDTD 'square rule' solver implemented with $\Delta$x=$\Delta$t grid, circular-buffer delay, analytical boundary conditions, and bilinear interpolation for $\chi$ construction. Figures 5, 6, 7 reproduced. Figure 8 (non-Markovian regime) not attempted --- requires \textasciitilde{}125 GB memory for full  history. Parallel performance benchmarking (paper's emphasis) not done; serial Python only.\par
\noindent\textbf{Agreement rationale.} Fig 5: zero boundary error, bounded interior error envelope, stable through 40K timesteps (confirmed). Fig 6 (ka=$\pi$/4): g(0)=0.001 at resonance (antibunching), 2.13 off-resonance (bunching) --- qualitatively exact. Fig 7 (ka=$\pi$/2): g(0)=0.007 at resonance, 0.539 off-resonance --- antibunching reproduced, bunching/oscillation features qualitatively present. No direct reference-solution comparison (paper's Ref [8] scattering theory not implemented).\par\smallskip
\noindent\textbf{What reproduced:}
\begin{itemize}[leftmargin=1.4em,itemsep=0.2em,topsep=0.2em]
\item Square-rule FDTD discretization (4-corner averages)
\item Stable  evolution through 40,000 timesteps (Fig 5)
\item Analytical boundary conditions from one-excitation sector e(t) with log-space bounce computation
\item $\chi$(x,x,t) construction with bilinear interpolation
\item g($\tau$) correlation function (Figs 6, 7)
\item Antibunching at resonance (g(0)$\rightarrow$0) and bunching off-resonance
\item Circular-buffer memory management for delay feedback
\end{itemize}\noindent\textbf{What missing:}
\begin{itemize}[leftmargin=1.4em,itemsep=0.2em,topsep=0.2em]
\item Figure 8 non-Markovian regime (requires \textasciitilde{}125 GB memory)
\item Parallel-performance benchmarks (paper's central contribution)
\item Reference scattering-theory solution for quantitative g comparison
\item Out-of-core or memory-mapped storage for large-grid runs
\item Numba/C JIT to realize paper-scale timing
\item Pipeline-parallel exploitation of causal delay structure
\item Multi-atom / multi-mirror generalizations
\end{itemize}\noindent\textbf{Follow-on research questions:}
\begin{enumerate}[leftmargin=1.6em,itemsep=0.2em,topsep=0.2em,label=Q\arabic*.]
\item Port the solver to CUDA and measure sustained throughput for (N, N) up to the Fig 8 regime --- at what grid does the GPU cross the memory wall, and does out-of-core NVMe backing restore feasibility
\item Implement the exact scattering-theory reference (Ref [8]) and quantify FDTD g error as a function of $\Delta$ --- does the empirical convergence rate match the paper's theoretical O($\Delta$$^2$)
\item Extend the PDE to N atoms at staggered positions --- does g(0) show monotonic deepening of antibunching, and is there a critical atom spacing where superradiance flips to subradiance
\item How does g(0) depend on ka beyond $\pi$/2 up to the Fig 8 regime (ka $\approx$ 10$\pi$--20$\pi$) Produce a phase diagram of g(0) vs (ka, $\omega$/) showing antibunching/bunching/chaotic zones.
\item Replace the perfect mirror with a finite-reflectivity boundary r\textless{}1 --- at what r does the non-Markovian oscillation structure collapse to Markovian exponential decay, and does an effective-Markovian approximation work beyond that threshold
\end{enumerate}

